\section{第13章\quad 无穷级数}
\subsection{常数项级数的概念与性质}
\begin{exercise}{655}{部分和与余项}
已知级数 $\sum_{n=1}^{\infty}u_n$ 的前 $n$ 项和为 $S_n=\frac{2n}{n+2}$。求 $u_1,u_2$ 及 $u_n$ 的表达式，按定义求级数的和 $S$，并求出使 $|S-S_n|<10^{-3}$ 的最小正整数 $n$。
\end{exercise}
\begin{exercise}{656}{裂项求和与误差}
设 $S_N=\sum_{n=1}^{N}\frac{3}{(3n-1)(3n+2)}$。求 $S_N$，据此证明相应无穷级数收敛并求其和；再求用 $S_N$ 近似该和时误差不超过 $10^{-3}$ 所需的最小 $N$。
\end{exercise}
\begin{exercise}{657}{对数级数的参数判别}
先求 $\sum_{n=1}^{N}\ln\left(1+\frac{2}{2n-1}\right)$ 并判定对应无穷级数的敛散性，再讨论实参数 $p$ 为何值时级数
\[\sum_{n=1}^{\infty}\frac{1}{n^p}\ln\left(1+\frac{2}{2n-1}\right)\]
收敛。说明临界参数处所用的判别依据。
\end{exercise}
\begin{exercise}{658}{定义与柯西准则}
对级数 $\sum_{n=0}^{\infty}(\sqrt{n+1}-\sqrt n)$，求前 $N+1$ 项和并判定敛散性；进一步计算 $\sum_{n=N}^{4N-1}(\sqrt{n+1}-\sqrt n)$，说明这一结果怎样反映该级数不满足柯西收敛准则。
\end{exercise}
\begin{exercise}{659}{根式裂项与绝对收敛}
对级数 $\sum_{n=1}^{\infty}(\sqrt[n+1]{2}-\sqrt[n]{2})$，求前 $N$ 项和并求无穷级数的和；利用通项符号判定其是否绝对收敛，再求使截断误差小于 $10^{-3}$ 的最小 $N$，可将答案保留为含对数的取整表达式。
\end{exercise}
\begin{exercise}{660}{三角数倒数与余项}
求 $\sum_{n=1}^{\infty}\frac1{1+2+\cdots+n}$ 的前 $N$ 项和及总和，给出精确余项；再讨论 $\sum_{n=1}^{\infty}\frac{n^p}{1+2+\cdots+n}$ 随实参数 $p$ 的敛散性。
\end{exercise}
\begin{exercise}{661}{有理化裂项与误差}
将 $\frac1{\sqrt{n+1}\sqrt n(\sqrt{n+1}+\sqrt n)}$ 写为两个相邻项之差，求相应级数的和及前 $N$ 项余项；求使误差不超过 $10^{-3}$ 的最小 $N$，并比较其与 $N^{-1/2}$ 的渐近关系。
\end{exercise}
\begin{exercise}{662}{等比级数的参数与和}
讨论实参数 $q$ 为何值时 $\sum_{n=1}^{\infty}(2-q)^n$ 收敛，求收敛时的和函数 $S(q)$；再求使 $S(q)=1$ 的全部 $q$，并验证所得参数属于收敛范围。
\end{exercise}
\begin{exercise}{663}{对数比值的几何级数}
先证明对任意 $a>1$ 都有 $0<\frac{\ln a}{a}<1$，再求 $S(a)=\sum_{n=1}^{\infty}(\frac{\ln a}{a})^n$，特别给出 $a=3$ 的值；求 $S(a)$ 在 $a>1$ 上的最大值及等号条件。
\end{exercise}
\begin{exercise}{664}{三角函数公比与临界角}
分别判定 $\sum_{n=0}^{\infty}\tan^n\theta$ 在 $\theta=\frac\pi6,\frac\pi4,\frac\pi3$ 时的敛散性，收敛时求和；再给出 $0<\theta<\frac\pi2$ 内的完整收敛范围，说明临界角处为何不能求和。
\end{exercise}
\begin{exercise}{665}{符号几何级数与绝对收敛}
对 $a>0$ 讨论 $\sum_{n=1}^{\infty}\frac{\cos n\pi}{a^n}$ 的敛散性与收敛类型，求收敛时的和；给出 $a=3$ 时的值及截取前 $N$ 项的精确余项。
\end{exercise}
\begin{exercise}{666}{两个幂次判据的交集}
求使 $\sum n^{-2p}$ 与 $\sum n^{-(3-p)}$ 同时收敛的全部实数 $p$，说明各临界点的情况；是否存在使两者同时发散的 $p$？证明结论。
\end{exercise}
\begin{exercise}{667}{必要条件与加括号的逆命题}
分别判断以下命题，成立者证明，不成立者给出通项明确的反例：
\begin{enumerate}
\item $\sum u_n$ 发散必推出 $u_n$ 不趋于零。
\item $\sum u_n,\sum v_n$ 均发散必推出 $\sum(2u_n-v_n)$ 发散。
\item $\sum u_n$ 发散必推出 $100+\sum_{n=1}^{\infty}u_{n+100}$ 发散。
\item $\sum u_n$ 发散必推出 $\sum_{n=1}^{\infty}(u_{2n-1}+u_{2n})$ 发散。
\end{enumerate}
对最后一项，补充 $u_n\to0$ 后结论是否改变？说明理由。
\end{exercise}
\begin{exercise}{668}{有限改变与逐项常数扰动}
已知 $\sum_{n=1}^{\infty}u_n=S$。逐一判定 $\sum1000u_n$、$\sum(u_n+\frac1{1000})$、$1000+\sum u_n$、$\sum u_{n+100}$ 的敛散性；收敛者用 $S$ 与必要的首项表示其和，发散者说明部分和趋向何处。
\end{exercise}
\begin{exercise}{669}{四类通项的判别依据}
分别判定下列级数的敛散性，给出各自最直接的判别依据，并求其中收敛级数的和：
\[\text{(1)}\ \sum_{n=1}^{\infty}\left(\frac1{2n}+\frac1{2^n}\right),\quad
\text{(2)}\ \sum_{n=1}^{\infty}\cos\frac1{n^2},\quad
\text{(3)}\ \sum_{n=1}^{\infty}\ln\left(1+\frac1n\right),\quad
\text{(4)}\ \sum_{n=1}^{\infty}(-1)^ne^{-n}.\]
对(3)另写出前 $N$ 项和，核验判定。
\end{exercise}
\begin{exercise}{670}{混合裂项的误差主项}
求 $\sum_{n=1}^{\infty}(3^{-n}+\frac3{(3n-1)(3n+2)})$ 的和及前 $N$ 项余项 $R_N$；求 $\lim_{N\to\infty}NR_N$，说明误差主要由哪一部分产生。
\end{exercise}
\begin{exercise}{671}{通项条件与参数}
讨论实参数 $a$ 为何值时级数
\[\sum_{n=1}^{\infty}\left(a\frac{n}{6n+1}+\frac{(-1)^n}{6^n}\right)\]
收敛，并在收敛时求和。对 $a=1$，再判断部分和趋于 $+\infty$、$-\infty$ 还是振荡，并说明理由。
\end{exercise}
\begin{exercise}{672}{级数线性运算}
已知实数项级数 $\sum_{n=1}^{\infty}u_n$ 收敛，$q\in\mathbb R\setminus\{0\}$。求使 $\sum_{n=1}^{\infty}(u_n+3q^{-n})$ 收敛的全部 $q$，并求两个级数的和之差。若原级数条件收敛，讨论新级数是否可能绝对收敛。
\end{exercise}
\begin{exercise}{673}{正项级数的临界指数}
讨论 $\sum_{n=1}^{\infty}(n^{1-p}+3\cdot2^{-n})$ 的敛散性，给出实参数 $p$ 的完整范围；分别说明在临界值处以及通项不趋于零的参数范围内，采用哪种判别方法即可完成判定。
\end{exercise}
\begin{exercise}{674}{线性组合与部分和}
已知 $\sum_{n=1}^{\infty}u_n=2$，且 $\sum_{k=1}^{n}v_k=\frac{2n}{3n+1}$。求 $v_n$ 及 $\sum_{n=1}^{\infty}(2u_n-3v_n)$；进一步求使 $\sum_{n=1}^{\infty}(a u_n-3v_n)$ 的和为零的实数 $a$，并说明逐项线性运算的依据。
\end{exercise}
\begin{exercise}{675}{奇偶分组与余项估计}
已知 $\sum_{n=1}^{\infty}n^{-2}=\frac{\pi^2}{6}$。求 $\sum_{n=1}^{\infty}(2n-1)^{-2}$ 及 $\sum_{n=1}^{\infty}(-1)^{n-1}n^{-2}$ 的和，并利用积分比较给出第一项级数截取前 $N$ 项后的余项上界。
\end{exercise}
\begin{exercise}{676}{收敛性质与反例}
设 $\sum u_n$ 收敛，$\sum v_n$ 发散。证明 $\sum(u_n+v_n)$ 必发散；分别构造满足上述条件且 $\sum u_nv_n$ 收敛、$\sum u_nv_n$ 发散的具体例子。各例均须写出通项并验证。
\end{exercise}
\begin{exercise}{677}{级数抵消与必要条件}
已知 $\sum(u_n+v_n)$ 收敛。证明 $\sum u_n$ 与 $\sum v_n$ 必同时收敛或同时发散；再构造一个 $u_n\to0,v_n\to0$ 且这两个级数均发散的例子，说明合并后的级数仍可收敛。
\end{exercise}
\begin{exercise}{678}{分组收敛与通项条件}
设 $\sum_{n=1}^{\infty}(u_{2n-1}+u_{2n})$ 收敛。证明原级数 $\sum u_n$ 收敛的充要条件为 $u_{2n}\to0$；举例说明若删去这一条件，分组级数可以收敛而原级数发散。
\end{exercise}
\begin{exercise}{679}{参数裂项与截断精度}
设 $a>0$，求 $\sum_{n=1}^{\infty}\frac{a}{(an+1)(a(n+1)+1)}$ 的和及前 $N$ 项余项，求使余项不超过 $\varepsilon>0$ 的最小正整数 $N$；讨论固定 $N$ 时余项随 $a$ 的变化。
\end{exercise}
\subsection{正项级数}
\begin{exercise}{680}{正项估计与拆分求和}
对 $\sum_{n=1}^{\infty}\frac{3+(-1)^{n+1}}{5^n}$，先给出两侧几何级数估计以证明收敛，再拆分求和；写出截取前 $N$ 项的精确余项，并验证它为正。
\end{exercise}
\begin{exercise}{681}{乘积项与几何上界}
利用几何级数比较判定下列级数的敛散性，并为各自从第 $N+1$ 项起的尾和给出明确上界：
\[\text{(1)}\ \sum_{n=1}^{\infty}\frac{3^{n-1}}{\prod_{k=1}^{n}(3k-2)},\qquad
\text{(2)}\ \sum_{n=1}^{\infty}\frac{2^n}{(2n-1)3^n}.\]
在(1)中明确给出相邻项比值及一个对所有 $n\ge1$ 成立的小于 $1$ 的上界。
\end{exercise}
\begin{exercise}{682}{不同增长阶的比较}
判定下列级数的敛散性，分别写出所选 $p$ 级数及通项比值的极限；对(3)进一步裂项求和：
\[\text{(1)}\ \sum_{n=1}^{\infty}\frac1{3n-1},\quad
\text{(2)}\ \sum_{n=1}^{\infty}\frac3{\sqrt n+\sqrt[3]n+\sqrt[4]n},\quad
\text{(3)}\ \sum_{n=1}^{\infty}\frac{n+3}{n(n+1)(n+2)}.\]
\end{exercise}
\begin{exercise}{683}{等价无穷小与尾和界}
利用极限比较判别下列级数，写出非零有限的比较极限：
\[\text{(1)}\ \sum_{n=1}^{\infty}\tan\frac\pi{3^n},\quad
\text{(2)}\ \sum_{n=1}^{\infty}\frac{3n^2+4}{n^4+1},\quad
\text{(3)}\ \sum_{n=1}^{\infty}\frac1{n\sqrt n}.\]
对(3)用积分比较给出前 $N$ 项余项的上下界，并说明误差随 $N$ 的阶。
\end{exercise}
\begin{exercise}{684}{三角函数中指数位置辨析}
分别判定 $\sum\sin\frac\pi6$、$\sum\sin\frac\pi{6n}$、$\sum\sin\frac\pi{6^n}$、$\sum\sin^n\frac\pi6$ 的敛散性，求第四个级数的和；对第三个级数给出几何级数尾和上界，解释这四种写法的通项极限或衰减速度。
\end{exercise}
\begin{exercise}{685}{对数等价与慢衰减}
判定 $\sum\ln(1+n^{-3})$、$\sum\ln(1+n^{-1/2})$、$\sum\frac1{\ln(3n+1)}$ 的敛散性，明确写出前两个的比较极限；再讨论 $\sum\ln(1+n^{-p})$ 随实参数 $p>0$ 的敛散性，处理临界值。
\end{exercise}
\begin{exercise}{686}{参数幂次与失效判据}
对 $p\in\mathbb R$ 讨论 $\sum_{n=1}^{\infty}n^{-p}\tan\frac1n$ 的敛散性，给出临界值处的比较极限；再求相邻项比值的极限，说明为什么比值判别法在本题不能区分参数范围。
\end{exercise}
\begin{exercise}{687}{已知上界能保证的收敛性}
设 $0\le u_n\le n^{-3}$。分别判定 $\sum u_n^3$、$\sum(u_n+n^{-3})$、$\sum(u_n+3^{-n})$、$\sum(u_n+\frac1{3n})$ 的敛散性；再求保证对所有满足条件的 $u_n$ 都有 $\sum n^p u_n$ 收敛的全部实数 $p$，并用具体例子证明范围不能扩大。
\end{exercise}
\begin{exercise}{688}{下界与渐近等价的作用}
设 $u_n>0$。
\begin{enumerate}
\item 若 $u_n\ge n^{-2}$，判断 $\sum nu_n$ 的敛散性，并说明能否据此判断 $\sum u_n$，不能时给出两种例子。
\item 若 $n^2u_n\to k>0$，判定 $\sum u_n$，并完整讨论 $\sum n^q u_n$ 随实参数 $q$ 的敛散性。
\end{enumerate}
\end{exercise}
\begin{exercise}{689}{平方改变收敛性的阈值}
设 $u_n=1-\cos\frac1{\sqrt n}$。求 $nu_n$ 的极限，判定 $\sum u_n$ 与 $\sum u_n^2$ 的敛散性；进一步讨论 $\sum u_n^p$ 对 $p>0$ 的收敛范围，说明取平方为什么改变判定。
\end{exercise}
\begin{exercise}{690}{阶乘与指数的比值判别}
分别用比值判别法处理 $\sum\frac1{(2n+1)!}$ 与 $\sum\frac{n!}{2^n}$，写出相邻项比值；对第二个级数，再将常数 $2$ 换为任意 $a>0$，说明是否可能通过选择 $a$ 使其收敛。
\end{exercise}
\begin{exercise}{691}{指数与阶乘的竞争}
判定 $\sum\frac{2^n}{n^2}$ 与 $\sum\frac{2^n n!}{n^n}$ 的敛散性，并求各自相邻项比值的极限；对收敛的级数，给出一个 $q<1$ 使充分大时尾项可由公比 $q$ 的几何级数控制。
\end{exercise}
\begin{exercise}{692}{奇数连乘与参数临界}
先用比值法判定 $\sum_{n=1}^{\infty}\frac{1\cdot3\cdots(2n-1)}{5^n n!}$；再将 $5$ 换为 $a>0$，求比值法能直接判定收敛或发散的参数范围。对临界值 $a=2$，利用
\[\frac{1\cdot3\cdots(2n-1)}{2^n n!}=\prod_{k=1}^{n}\left(1-\frac1{2k}\right)\ge\frac1{n+1}\]
判定敛散性，并以数学归纳法证明所用不等式。
\end{exercise}
\begin{exercise}{693}{指数衰减与和函数求和}
证明 $\sum_{n=1}^{\infty}\frac1{(n+1)e^n}$ 收敛；由 $-\ln(1-x)=\sum_{n=1}^{\infty}x^n/n$ 求其和，并给出用几何级数比较得到的一个前 $N$ 项余项上界。
\end{exercise}
\begin{exercise}{694}{两个指数的渐近比较}
对 $a>0$ 讨论 $\sum_{n=1}^{\infty}\frac{a^n}{7^n-2^n}$ 的敛散性，特别给出 $a=5$ 时的判定；说明在临界参数处为何必须检验通项，不能只引用比值极限为 $1$。
\end{exercise}
\begin{exercise}{695}{含参数系数的几何型级数}
设 $a>0$，讨论 $\sum_{n=1}^{\infty}\frac{an+1}{a^n}$ 的敛散性；在收敛范围内利用几何级数及其导数求和，并单独解释 $a=1$ 时通项的行为。
\end{exercise}
\begin{exercise}{696}{周期系数的拆分求和}
证明 $\sum_{n=1}^{\infty}\frac{3n+\sin^2(n\pi/3)}{3^n}$ 收敛；利用 $\sin^2(n\pi/3)$ 的周期性拆分为几何级数，求总和，并说明分组合法的依据。
\end{exercise}
\begin{exercise}{697}{递推通项与零初值例外}
设 $0\le a_n\le b_n$，且 $a_{n+1}=\frac12(1+\frac1n)^n a_n$。当 $a_1>0$ 时，判定 $\sum a_n$ 与 $\sum b_n$ 的敛散性；当 $a_1=0$ 时，说明是否仍可唯一判断 $\sum b_n$，并给出必要的反例。
\end{exercise}
\begin{exercise}{698}{混合级数与尾和估计}
对从 $n=1$ 开始的级数 $\sum_{n=1}^{\infty}(n^{-4}+\frac{4^n}{(n+1)!})$，分别判定两部分并证明总级数收敛；当 $N\ge7$ 时，利用后一部分相邻项比值不超过 $1/2$ 及积分比较，为总余项给出明确上界。
\end{exercise}
\begin{exercise}{699}{用级数收敛推出通项极限}
令 $u_n=\frac{(n!)^2}{(2n+1)!}$。先证明 $\sum_{n=1}^{\infty}u_n$ 收敛，由收敛必要条件求 $\lim u_n$；再找一个对所有 $n\ge1$ 成立的几何上界，以定量说明该极限。
\end{exercise}
\begin{exercise}{700}{根值法与统一几何控制}
对 $\sum_{n=1}^{\infty}(\frac{2n}{4n+1})^n$ 求根值极限并判定敛散性；进一步给出一个对全部 $n$ 成立的几何级数上界，并据此估计从第 $N+1$ 项起的尾和。
\end{exercise}
\begin{exercise}{701}{指数项的参数阈值}
讨论 $\sum_{n=1}^{\infty}\frac{a^n}{2^n(3n+1)}$ 随 $a>0$ 的敛散性，给出 $a=3$ 时的判定；对临界参数用比较法处理，并说明分母中的一次多项式为何不能压过指数增长。
\end{exercise}
\begin{exercise}{702}{根值判别与比较估计}
对下列两个级数，分别求通项的 $n$ 次方根的极限并判定敛散性。对收敛的级数，给出一个 $0<q<1$ 及整数 $N$，使其第 $n$ 项在 $n\ge N$ 时不超过 $q^n$。
\[\text{(1)}\quad\sum_{n=1}^{\infty}\left(\frac{4n-3}{3n+2}\right)^{2n};\qquad
\text{(2)}\quad\sum_{n=1}^{\infty}\left(1-\frac1n\right)^{n^2}.\]
\end{exercise}
\begin{exercise}{703}{指数项与多项式因子}
判定下列级数的敛散性，并说明乘上 $n^{-2}$ 这样的因子是否一定能使指数增长的通项趋于零。
\[\text{(1)}\quad\sum_{n=1}^{\infty}\frac{(\arctan n)^n}{2^n};\qquad
\text{(2)}\quad\sum_{n=1}^{\infty}\frac1{n^2}\left(1+\frac1n\right)^{n^2}.\]
对(1)给出一个可用于比较判别的几何级数。
\end{exercise}
\begin{exercise}{704}{根值判别的临界情形}
设 $a>0,p\in\mathbb R$。先判定 $\sum_{n=1}^{\infty}(\frac{na}{n+1})^n$ 的敛散性，再讨论
\[\sum_{n=1}^{\infty}\frac1{n^p}\left(\frac{na}{n+1}\right)^n\]
的敛散性。须单独处理根值判别法不能直接判定的参数值。
\end{exercise}
\begin{exercise}{705}{判别法失效与反例}
分别构造正项级数，使其相邻项比值极限均为 $1$，而一个收敛、一个发散；再证明若 $u_n>0$ 且 $n^2u_n\to L\in(0,\infty)$，则 $\sum u_n$ 收敛。说明极限比较与比值判别法所用的信息有何区别。
\end{exercise}
\begin{exercise}{706}{相消后的参数阈值}
对实参数 $p$ 讨论正项级数 $\sum_{n=1}^{\infty}n^p\left(\frac1n-\ln(1+\frac1n)\right)$ 的敛散性。先用定积分或二阶泰勒公式求括号内量与 $n^{-2}$ 的比值极限，证明正项性，并单独处理临界参数。
\end{exercise}
\subsection{任意项级数}
\begin{exercise}{707}{交错条件与绝对值比较}
分别判定 $\sum_{n=1}^{\infty}\frac{(-1)^n}{(2n)!}$ 与 $\sum_{n=1}^{\infty}\frac{(-1)^n}{\sqrt{n^2+a}+1}$（$a>0$）的收敛类型；对两者给出莱布尼茨余项界，说明参数 $a$ 是否改变第二个级数的收敛类型。
\end{exercise}
\begin{exercise}{708}{交错判别与通项条件}
对下列级数分别判定绝对收敛、条件收敛或发散，并指出莱布尼茨判别法的条件是否满足。
\[\text{(1)}\quad\sum_{n=1}^{\infty}(-1)^n\sin\frac1{\sqrt n};\qquad
\text{(2)}\quad\sum_{n=1}^{\infty}(-1)^n\cos\frac1{\sqrt n}.\]
对收敛的级数，估计取前 $N$ 项时的误差，并给出误差小于 $10^{-2}$ 的一个充分条件。
\end{exercise}
\begin{exercise}{709}{交错级数与绝对收敛}
判定下列级数的收敛类型，并明确验证绝对值通项的单调性及极限：
\[\text{(1)}\quad\sum_{n=1}^{\infty}(-1)^n\ln\left(1+\frac1n\right);\qquad
\text{(2)}\quad\sum_{n=2}^{\infty}\frac{(-1)^n}{\ln n}.\]
比较用莱布尼茨余项估计使误差小于 $10^{-3}$ 时，两者所需项数的数量级。
\end{exercise}
\begin{exercise}{710}{绝对收敛与符号拆分}
证明 $\sum_{n=1}^{\infty}(-1)^n\frac{2+(-1)^n}{4^n}$ 绝对收敛；把它拆为两个几何级数求和，并写出前 $N$ 项余项。说明本题按奇偶项分组求和的合法性。
\end{exercise}
\begin{exercise}{711}{交错形式与单调性缺失}
按所列顺序讨论级数 $\frac1{\sqrt2-1}-\frac1{\sqrt2+1}+\frac1{\sqrt3-1}-\frac1{\sqrt3+1}+\cdots$ 的敛散性。写出前 $2N$ 项和，说明通项虽趋于零且符号交替，为何仍不能用莱布尼茨判别法推出收敛；再判断奇数项部分和的极限行为。
\end{exercise}
\begin{exercise}{712}{幂指数与交错收敛类型}
分别判定 $\sum n^{-1/2}$、$\sum(-1)^nn^{-1/2}$、$\sum n^{-3/2}$、$\sum(-1)^nn^{-3/2}$ 的敛散性与收敛类型。证明对一般 $p>0$，正项级数 $\sum n^{-p}$ 与交错级数 $\sum(-1)^nn^{-p}$ 的收敛条件分别是什么，明确两个临界点的作用。
\end{exercise}
\begin{exercise}{713}{周期符号与稀疏交错项}
判定 $\sum_{n=1}^{\infty}\frac1n\sin\frac{n\pi}2$ 的收敛类型，写成去掉零项后的级数并求和；给出截取前 $4N$ 项的误差上界，说明插入零项为何不影响和。
\end{exercise}
\begin{exercise}{714}{交错判别的单调性证明}
判定 $\sum_{n=1}^{\infty}\frac{(-1)^n}{n-\ln n}$ 的收敛类型。证明分母始终为正、绝对值通项最终单调趋零，并给出前 $N$ 项误差界；绝对收敛部分须写出与调和级数的比较极限。
\end{exercise}
\begin{exercise}{715}{柯西不等式与最优上界}
设实数列满足 $\sum u_n^2<\infty$。证明 $\sum\frac{u_n}{n}$ 绝对收敛，并证明 $\sum\frac{|u_n|}{n}\le\frac\pi{\sqrt6}(\sum u_n^2)^{1/2}$；给出非零数列说明常数 $\pi/\sqrt6$ 不能减小。
\end{exercise}
\begin{exercise}{716}{平移交错级数的统一误差}
设 $x>0$，判定 $\sum_{n=1}^{\infty}\frac{(-1)^n}{n+x}$ 的收敛类型；证明对所有 $x>0$，前 $N$ 项误差均不超过 $1/(N+1)$，并据此给出不依赖 $x$ 的截断项数，使误差小于 $10^{-3}$。
\end{exercise}
\begin{exercise}{717}{有界振荡与几何控制}
判定 $\sum_{n=1}^{\infty}\frac{(-1)^n}{\pi^{n+1}}\cos\frac{n\pi}{n+1}$ 的收敛类型，给出绝对值级数的几何上界；再估计从第 $N+1$ 项起的尾和绝对值，说明无需检验余弦因子的单调性。
\end{exercise}
\begin{exercise}{718}{参数谐波项的分离}
求使 $\sum_{n=1}^{\infty}\frac{(-1)^n+kn}{n^2}$ 收敛的全部实数 $k$，说明收敛时的类型并求和，可用 $\sum n^{-2}=\pi^2/6$；对其余 $k$，判断部分和趋于何处。
\end{exercise}
\begin{exercise}{719}{四个交错级数的分类}
分别判定 $\sum(-1)^{n-1}n^{-2/3}$、$\sum(-1)^{n-1}n^{-3/2}$、$\sum(-1)^{n-1}(2/3)^n$、$\sum(-1)^{n-1}(3/2)^n$ 的收敛类型；对条件收敛者给出截断误差界，对几何型收敛者求和。
\end{exercise}
\begin{exercise}{720}{有界系数与绝对收敛}
分别判定下列级数的收敛类型，须给出绝对值比较或通项不趋零的依据：
\[\text{(1)}\ \sum\frac{(-1)^n(3+(-1)^n)}{(3n+1)\sqrt{n+3}},\qquad
\text{(2)}\ \sum\frac{(-1)^n2^n}{n^2+1},\]
\[\text{(3)}\ \sum\frac{(-1)^n\sqrt{n^2+1}}{n+1+\sin n},\qquad
\text{(4)}\ \sum\frac{(-1)^n}{3n-1}.\]
对(1)进一步说明若把分子中的 $3+(-1)^n$ 换成任意有界实数列，是否仍绝对收敛。
\end{exercise}
\begin{exercise}{721}{交错符号并非充分条件}
分别分类以下四个级数，并为每个给出关键极限或单调性论证：
\[\text{(1)}\ \sum(-1)^n\left(\frac n{n+1}\right)^n,\quad
\text{(2)}\ \sum(-1)^n3^n\sin\frac\pi{5^n},\]
\[\text{(3)}\ \sum(-1)^n\frac{n+2}{(n+1)\sqrt n},\quad
\text{(4)}\ \sum(-1)^n\frac{5^n-3^n}{5^n}.\]
对(3)给出前 $N$ 项的莱布尼茨误差界。
\end{exercise}
\begin{exercise}{722}{振荡系数与绝对值比值}
证明对每个 $x>0$，$\sum_{n=1}^{\infty}\frac{n^n}{3^nn!}\cos nx$ 均绝对收敛；对不含余弦因子的正项级数计算相邻项比值极限，并用同一尾和界同时控制全部 $x>0$ 时的截断误差。
\end{exercise}
\begin{exercise}{723}{平方可和与绝对收敛}
设 $k>0$，实数项级数 $\sum_{n=1}^{\infty}u_n^2$ 收敛。利用柯西不等式证明 $\sum_{n=1}^{\infty}\frac{u_n}{\sqrt{n^2+k}}$ 绝对收敛，并给出其绝对值级数的和关于 $\sum u_n^2$ 的一个上界。讨论把 $k>0$ 改为 $k=0$ 后结论是否仍成立。
\end{exercise}
\begin{exercise}{724}{交错级数与平方级数}
设 $u_n=(-1)^n\tan\frac1{\sqrt n}$。判定 $\sum u_n$ 的收敛类型及 $\sum u_n^2$ 的敛散性；进一步讨论正项级数 $\sum |u_n|^p$ 随实参数 $p>0$ 变化的敛散性，并处理临界值。
\end{exercise}
\begin{exercise}{725}{条件收敛与参数分解}
对实参数 $k$，讨论 $\sum_{n=1}^{\infty}\frac{k+n}{n^2}\cos n\pi$ 的收敛类型，并利用 $\sum n^{-2}=\frac{\pi^2}{6}$ 与 $\sum_{n=1}^{\infty}(-1)^{n-1}/n=\ln2$ 求其和。说明改变有限多个通项的符号是否影响收敛类型。
\end{exercise}
\begin{exercise}{726}{级数夹逼的证明}
设实数列满足 $u_n\le a_n\le v_n$，且 $\sum u_n,\sum v_n$ 都收敛。以非负级数 $\sum(a_n-u_n)$ 为出发点证明 $\sum a_n$ 收敛；再判断能否推出 $\sum|a_n|$ 收敛，若不能，给出满足全部条件的反例。
\end{exercise}
\begin{exercise}{727}{交错收敛的条件辨析}
设 $a_n>0,a_n\to0$。判断 $\sum(-1)^na_n$ 是否必收敛；构造明确数列使其发散，并指出该数列不满足的莱布尼茨条件。再证明若 $a_n$ 最终单调递减，则改变有限个首项不影响收敛性。
\end{exercise}
\begin{exercise}{728}{带正项扰动的交错级数}
对实参数 $p>0,q>0$ 讨论 $\sum_{n=1}^{\infty}((-1)^{n-1}n^{-p}+n^{-q})$ 的收敛、绝对收敛或条件收敛范围。须区分 $p=q$ 与 $p\ne q$，并证明绝对值级数的临界情况。
\end{exercise}
\subsection{幂级数}
\begin{exercise}{729}{收敛半径与边界通项}
求 $\sum_{n=0}^{\infty}\frac{(-1)^n(3n+1)}{3^{n+1}}x^n$ 的收敛半径及收敛域；在收敛区间内拆成几何级数及其导数求和，并说明两个边界点为何不能直接代入所得有理函数作为级数和。
\end{exercise}
\begin{exercise}{730}{两种指数系数的主导项}
求 $\sum_{n=1}^{\infty}\frac{3^n+(-2)^n}{n^2}x^n$ 的收敛半径、收敛域及端点处的收敛类型；再将 $n^2$ 改为 $n^p$（$p>0$），讨论两个端点的收敛性随 $p$ 的变化。
\end{exercise}
\begin{exercise}{731}{稀疏幂次的半径}
设 $a_n$ 最终非零且 $|a_{n+1}/a_n|\to4$。求 $\sum_{n=0}^{\infty}a_nx^{2n+1}$ 的收敛半径；说明能否仅凭该条件确定两个端点的收敛性，并分别给出端点都绝对收敛和端点都发散的系数例子。
\end{exercise}
\begin{exercise}{732}{条件收敛点决定半径}
已知 $\sum_{n=1}^{\infty}a_n(x+1)^n$ 在 $x=-2$ 条件收敛。证明其收敛半径等于 $1$，判断能否确定另一个端点的敛散性；若不能，分别给出另一个端点收敛与发散的系数例子，并写出相应完整收敛域。
\end{exercise}
\begin{exercise}{733}{阶乘系数的整函数级数}
求 $\sum_{n=1}^{\infty}\frac{n^2}{(2n)!}x^n$ 的收敛域；当 $x>0$ 时，令 $t=\sqrt x$，从 $\cosh t=\sum_{n=0}^{\infty}t^{2n}/(2n)!$ 出发求其和函数，并求在 $x=0$ 的连续取值。
\end{exercise}
\begin{exercise}{734}{临界端点与通项增长}
求 $\sum_{n=1}^{\infty}\frac{n!}{n^n}x^n$ 的收敛区间及收敛域。处理端点时，可利用 $(1+1/n)^n<e$ 证明 $e^nn!/n^n$ 随 $n$ 增加，进而检验通项是否趋零；无须使用斯特林公式。
\end{exercise}
\begin{exercise}{735}{端点收敛类型与求导}
求 $\sum_{n=1}^{\infty}\frac{(-1)^n}{2^n\sqrt n}x^n$ 的收敛域，分别判断端点的收敛类型；再求逐项求导后的幂级数的收敛域，说明求导保持半径是否也保持端点性质。
\end{exercise}
\begin{exercise}{736}{平移中心与导数级数}
求 $\sum_{n=1}^{\infty}\frac{(x-2)^n}{n^2+1}$ 的收敛域及端点收敛类型；求其导数级数的收敛半径，分别判定导数级数在两个端点的敛散性。
\end{exercise}
\begin{exercise}{737}{内点绝对收敛的定量证明}
已知 $\sum a_nx^n$ 在 $x=-3$ 收敛。证明它在 $x=2$ 绝对收敛；进一步，对任意 $0<r<3$，用 $|a_n(-3)^n|\le M$ 给出 $|x|\le r$ 上从第 $N+1$ 项起的绝对值尾和上界。
\end{exercise}
\begin{exercise}{738}{中心与距离判定}
设 $\sum_{n=1}^{\infty}a_n(x-1)^n$ 的收敛半径为 $1$。判定它在 $x=3$ 的敛散性，并说明在 $x=0,2$ 处能否确定；分别给出半径恰为 $1$ 且端点均收敛、均发散的系数例子。
\end{exercise}
\begin{exercise}{739}{两个等距端点确定收敛域}
已知 $\sum_{n=1}^{\infty}a_n(x+\frac12)^n$ 在 $x=1$ 收敛、在 $x=-2$ 发散。求收敛半径及完整收敛域，证明 $x=1$ 处必为条件收敛；给出满足要求的一个系数序列。
\end{exercise}
\begin{exercise}{740}{远处点的不可确定性}
已知 $\sum_{n=1}^{\infty}a_n(x-2)^n$ 在 $x=-1$ 收敛。说明能否唯一判断 $x=6$ 处的收敛性；分别构造在 $x=6$ 绝对收敛、条件收敛、发散的系数，并核验每个例子在 $x=-1$ 的条件。
\end{exercise}
\begin{exercise}{741}{逐项求导与端点}
已知 $\sum_{n=0}^{\infty}a_nx^n$ 的收敛半径为 $3$。求 $\sum_{n=1}^{\infty}na_n(x-2)^n$ 的收敛半径与收敛区间；说明能否仅凭已知条件判断其在 $x=-1,5$ 处的收敛性，并分别举例说明端点收敛与发散均可能发生。
\end{exercise}
\begin{exercise}{742}{几何级数与逐项积分}
求 $S(x)=\sum_{n=1}^{\infty}\frac{x^{n+1}}{3^n}$ 的收敛域及和函数，并说明在任意 $0<r<3$ 上逐项积分成立的理由，进而计算 $\int_0^r S(x)\,\dd x$。
\end{exercise}
\begin{exercise}{743}{望远镜求和与一致误差}
求 $S_N(x)=\sum_{n=1}^{N}(\frac{x^n}{n}-\frac{x^{n+1}}{n+1})$。据此求该级数在 $[-1,1]$ 上的和函数，并直接估计 $\sup_{|x|\le1}|S(x)-S_N(x)|$，判断是否一致收敛。
\end{exercise}
\begin{exercise}{744}{对数和函数与边界}
求 $\sum_{n=1}^{\infty}\frac{5^n}{n}x^n$ 的收敛域及和函数；单独判定两个边界点的收敛类型，并求在 $x=-\frac15$ 处级数的和，说明该值怎样与区间内部的和函数衔接。
\end{exercise}
\begin{exercise}{745}{幂级数求和与积分}
求 $S(x)=\sum_{n=1}^{\infty}\frac{(-1)^{n-1}}n x^{n+1}$ 的收敛域及和函数，分别处理两个端点；再利用和函数计算 $\sum_{n=1}^{\infty}\frac{(-1)^{n-1}}{n(n+2)}$，说明逐项积分所采用的区间与依据。
\end{exercise}
\begin{exercise}{746}{反正切级数与可去奇点}
求 $\sum_{n=0}^{\infty}\frac{(-1)^n}{2n+1}x^{2n}$ 的收敛域及和函数 $S(x)$，特别写明 $S(0)$；求 $S''(0)$，并利用端点处级数得到 $\sum_{n=0}^{\infty}\frac{(-1)^n}{2n+1}$ 的和。
\end{exercise}
\begin{exercise}{747}{逐项求导与积分应用}
从几何级数出发求 $S(x)=\sum_{n=1}^{\infty}nx^{2n-1}$ 的收敛域及和函数；求 $\int_0^{1/2}S(x)\,\dd x$，并说明为何可以在该积分区间内逐项积分。
\end{exercise}
\begin{exercise}{748}{展开中心与和函数}
求 $S(x)=\sum_{n=1}^{\infty}n(x+2)^n$ 的收敛域及和函数；在收敛区间内求 $S$ 的零点与单调区间，并判定级数在 $x=-3,-1$ 处是否收敛。
\end{exercise}
\begin{exercise}{749}{求导保持半径但改变边界}
对 $F(x)=\sum_{n=1}^{\infty}x^n/n^2$，分别求原级数、一次导数级数、二次导数级数的收敛域，证明三者收敛半径相同；说明为何不能把开区间内的逐项求导结论无条件推广到端点。
\end{exercise}
\begin{exercise}{750}{参数幂级数与可导区间}
设 $p>0$，$F_p(x)=\sum_{n=1}^{\infty}(x-1)^n/n^p$。求其收敛域及导数级数的收敛域，完整讨论两个端点随 $p$ 的变化；对任意 $0<r<1$，给出 $|x-1|\le r$ 上的几何控制，说明和函数在开区间内连续可导。
\end{exercise}
\subsection{函数展开成幂级数}
\begin{exercise}{751}{复合指数函数的系数}
将 $e^{-x^2}$ 展开成 $x$ 的幂级数并求收敛域，求其在原点的一般奇数阶与偶数阶导数；用展开至 $x^4$ 的多项式近似 $e^{-1/4}$，写出交错级数误差上界。
\end{exercise}
\begin{exercise}{752}{正弦展开与一般系数}
求 $\sin2x$ 的麦克劳林展开及 $x^7$ 的系数；对 $a\in\mathbb R$ 写出 $\sin(ax)$ 中 $x^{2m+1}$ 的系数，说明偶数次项为何全部为零，并给出展开的收敛域。
\end{exercise}
\begin{exercise}{753}{几何展开的常数项核验}
将 $1/(3-x)$ 展开成 $x$ 的幂级数，写出通项系数、收敛域及截断到 $x^N$ 的精确余项；通过 $x=0$ 的取值说明为何不能遗漏首项或把系数写成 $3^{-n}$。
\end{exercise}
\begin{exercise}{754}{三角恒等式与幂级数展开}
利用倍角公式将 $\sin^2x$ 展开成 $x$ 的幂级数，写出前四个非零项、通项及收敛域；求该函数在原点的 $2m$ 阶导数，并解释其奇数阶导数的取值。
\end{exercise}
\begin{exercise}{755}{平方根代换与实变量延拓}
求 $S(x)=\sum_{n=0}^{\infty}\frac{(-1)^nx^{n+1}}{(2n+1)!}$ 的收敛域。分别在 $x>0,x=0,x<0$ 写出实值和函数，可使用 $\sinh t=(e^t-e^{-t})/2$；说明原级数为何消除了平方根表达式在原点的表面限制。
\end{exercise}
\begin{exercise}{756}{逐项积分与端点求和}
从 $1/(1+x)=\sum_{n=0}^{\infty}(-1)^nx^n$ 出发推导 $\ln(1+x)$ 的展开，求完整收敛域；说明在 $x=1$ 与 $x=-1$ 两端的情况，再利用展开给出 $\ln\frac32$ 的三项近似及误差界。
\end{exercise}
\begin{exercise}{757}{泰勒展开与交错余项}
从 $\frac1{1+x^2}$ 的几何级数出发，推导 $\arctan x$ 的麦克劳林展开并给出收敛域；用前三个非零项近似 $\arctan\frac12$，写出严格误差上界。
\end{exercise}
\begin{exercise}{758}{有理函数的幂级数}
将 $f(x)=\frac1{(3-x)^2}$ 展开成 $x$ 的幂级数，给出通项系数及收敛域；利用系数求 $f^{(m)}(0)$，其中 $m$ 为非负整数，并说明两个端点均不在收敛域内的理由。
\end{exercise}
\begin{exercise}{759}{部分分式与收敛半径}
将 $f(x)=\frac1{x^2-6x+5}$ 展开成 $x$ 的幂级数，写出通项系数并求收敛域；说明半径为何由距离展开中心最近的实奇点决定，再求 $f^{(6)}(0)$。
\end{exercise}
\begin{exercise}{760}{平移中心的指数展开}
将 $e^{3x}$ 展开成 $(x-1)$ 的幂级数，说明其对全部实数成立；用该展开构造 $e^{3x}$ 在 $x=1$ 附近的三次泰勒多项式，并用拉格朗日余项估计其在 $|x-1|\le\frac1{10}$ 上的最大误差。
\end{exercise}
\begin{exercise}{761}{平移几何级数与余项}
将 $f(x)=\frac1{3-x}$ 展开成 $(x+1)$ 的幂级数并求收敛域；写出截断至 $(x+1)^N$ 的精确余项，并据此给出 $|x+1|\le2$ 上的统一误差上界。
\end{exercise}
\begin{exercise}{762}{对数的平移展开与端点}
利用 $\ln(1+t)$ 的展开推导 $\ln x$ 关于 $(x-2)$ 的幂级数，并求收敛域；说明 $x=0,4$ 两个端点的不同情形，再用展开的前两项近似 $\ln\frac52$ 并估计误差。
\end{exercise}
\begin{exercise}{763}{负中心展开与导数}
将 $x^{-2}$ 展开成 $(x+1)$ 的幂级数，明确收敛域及通项系数；由此求 $f^{(n)}(-1)$，并说明将 $x=-2$ 或 $x=0$ 代入该级数时为何不能得到有限的级数和。
\end{exercise}
\begin{exercise}{764}{部分分式与中心平移}
将 $f(x)=\frac{x}{x^2-3x+2}$ 展开成 $(x+2)$ 的幂级数，求其收敛域及通项系数；利用所得系数求 $f^{(4)}(-2)$，并单独验证收敛区间的两个端点。
\end{exercise}
\begin{exercise}{765}{泰勒系数与函数相等的条件}
设 $f$ 在原点邻域具有各阶导数。说明写出形式泰勒级数是否足以证明它等于 $f$。对 $f(x)=e^x$，用拉格朗日余项证明其麦克劳林展开在任意固定闭区间 $[-M,M]$ 上成立，并给出余项的统一上界。
\end{exercise}
\begin{exercise}{766}{有理函数的展开中心选择}
将 $f(x)=1/((1-x)(2-x))$ 分别展开为 $x$ 的幂级数与 $(x+1)$ 的幂级数，写出通项系数和收敛域；求两展开在 $x=0$ 的取值，说明展开中心改变为何会改变收敛区间却不改变函数值。
\end{exercise}
\subsection{幂级数的应用}
\begin{exercise}{767}{系数提取与一般导数}
设 $f(x)=2x^2/(1+x)$。用幂级数求 $f^{(5)}(0)$，再写出 $f^{(n)}(0)$ 对全部非负整数 $n$ 的表达式，分别处理 $n=0,1$；说明所用展开在哪个邻域成立。
\end{exercise}
\begin{exercise}{768}{乘积展开的奇偶阶导数}
设 $f(x)=(1+x)\arctan x$。求 $f^{(100)}(0),f^{(101)}(0)$，并给出一般偶数阶与奇数阶导数的公式；写明幂级数乘法中各阶系数来自哪些项。
\end{exercise}
\begin{exercise}{769}{指数乘积的两种求导方法}
设 $f(x)=xe^{-x}$。利用麦克劳林系数求 $f^{(n)}(0)$（$n\ge1$），再用莱布尼茨乘积求导公式核验；求 $\sum_{n=1}^{\infty}f^{(n)}(0)x^n/n!$ 的收敛域与和函数。
\end{exercise}
\begin{exercise}{770}{平移对数的高阶导数}
对 $f(x)=\ln(x^2+4x+5)$，写出在 $x=-2$ 附近的幂级数，求全部 $n\ge1$ 的 $f^{(n)}(-2)$；说明奇数阶导数为零的对称性依据，并给出展开的收敛域。
\end{exercise}
\begin{exercise}{771}{分式展开与导数恢复}
将 $f(x)=(x-1)/(2x+1)$ 展开成 $(x-1)$ 的幂级数，求其收敛域及 $f^{(n)}(1)$（$n\ge1$）；再写出截断至 $(x-1)^N$ 的精确余项，说明最接近展开中心的奇点在哪里。
\end{exercise}
\begin{exercise}{772}{积分几何级数求数项和}
构造 $F(x)=\sum_{n=1}^{\infty}x^n/(n+1)$，先求其收敛域与和函数，再求 $\sum_{n=1}^{\infty}1/(2^n(n+1))$ 的和；在 $x=0$ 单独写出和函数的连续值。
\end{exercise}
\begin{exercise}{773}{线性加权几何级数}
求 $F(x)=\sum_{n=0}^{\infty}(n+1)x^n$ 的收敛域及和函数，进而求 $\sum_{n=0}^{\infty}(n+1)/3^n$；通过对有限几何和求导，给出后一数项级数前 $N+1$ 项的精确余项。
\end{exercise}
\begin{exercise}{774}{积分求和与端点极限}
设 $F(x)=\sum_{n=0}^{\infty}\frac{(-1)^n x^{3n+1}}{3n+1}$。先在 $|x|<1$ 内求 $F'(x)$ 并表示 $F(x)$；再说明能否令 $x\to1^-$ 求出 $\sum_{n=0}^{\infty}\frac{(-1)^n}{3n+1}$，并将其和写成初等常数。
\end{exercise}
\begin{exercise}{775}{三角级数与加权求和}
从 $\cos x$ 的麦克劳林展开出发，求 $F(x)=\sum_{n=0}^{\infty}\frac{(-1)^n(2n+1)}{(2n)!}x^{2n}$ 的和函数及收敛域，进而求 $F(1)$；说明使用逐项求导的依据。
\end{exercise}
\begin{exercise}{776}{二次加权幂级数}
求 $S(x)=\sum_{n=1}^{\infty}\frac{n(n+1)}{2^n}x^n$ 的收敛域及和函数，进而求 $\sum_{n=1}^{\infty}\frac{(-1)^n n(n+1)}{2^n}$ 的和；再用同一和函数求 $\sum_{n=1}^{\infty}\frac{n(n+1)}{4^n}$。
\end{exercise}
\begin{exercise}{777}{端点代入与逐项求导辨析}
对 $F(x)=\sum_{n=1}^{\infty}(-1)^{n-1}x^n/n$，求收敛域与和函数。判断能否在 $x=1$ 直接把导数级数求和得到 $F$ 在该点的左导数；写出导数级数的通项行为，并说明它与 $\ln(1+x)$ 在该点可导之间为何不矛盾。
\end{exercise}
\begin{exercise}{778}{高阶系数与交错数项和}
设 $f(x)=\ln(1+x^2)$。求全部 $f^{(n)}(0)$，再构造幂级数求 $\sum_{n=1}^{\infty}(-1)^{n-1}/(n2^n)$ 及 $\sum_{n=1}^{\infty}(-1)^{n-1}/(n(n+1)2^{n+1})$ 的和。明确逐项积分区间及合法性。
\end{exercise}
\subsection{傅里叶级数}
\begin{exercise}{779}{非对称周期段与常数项}
设 $f$ 为 $2\pi$ 周期函数，在 $0<x\le2\pi$ 上 $f(x)=x^2$。求傅里叶系数 $a_0$ 与 $a_n$（$n\ge1$），并求傅里叶级数在 $x=0,2\pi$ 处的和；说明积分区间可否从 $[0,2\pi]$ 改为任意长度 $2\pi$ 的区间。
\end{exercise}
\begin{exercise}{780}{周期归约与跳跃平均}
将 $f(x)=x^3+1$（$-\pi\le x<\pi$）作 $2\pi$ 周期延拓。求其傅里叶级数在 $x=7\pi/2,\pi,0$ 处的和，并写出全体间断点处的和；解释在间断点改变 $f$ 的单点取值为何不改变傅里叶系数。
\end{exercise}
\begin{exercise}{781}{分段周期函数的和函数}
设 $f$ 为 $2\pi$ 周期函数，在 $[-\pi,\pi)$ 上为 $f(x)=x$（$x<0$），$f(x)=0$（$x\ge0$）。求其傅里叶级数的和函数在 $[-\pi,\pi]$ 上的表达式，给出 $S(3\pi)$，并求常数项 $a_0/2$。
\end{exercise}
\begin{exercise}{782}{正弦展开的奇延拓}
设 $f(x)=|x-\pi/2|$（$0\le x\le\pi$），$b_n=\frac2\pi\int_0^\pi f(x)\sin(nx)\,\dd x$。记 $S(x)=\sum_{n=1}^{\infty}b_n\sin(nx)$。说明对应奇延拓的周期与间断点，求 $S(-9\pi/4),S(0),S(\pi)$；解释端点处级数和与原函数端点值的区别。
\end{exercise}
\begin{exercise}{783}{分段线性函数的完整展开}
将 $f(x)=2x$（$-\pi\le x<0$）、$f(x)=4x$（$0\le x<\pi$）作 $2\pi$ 周期延拓。先分解为奇函数与偶函数之和，求全部傅里叶系数及级数的和函数 $S(x)$；明确 $x=(2k+1)\pi$ 与 $x=2k\pi$ 处的和，$k\in\mathbb Z$。
\end{exercise}
\begin{exercise}{784}{周期延拓与间断点取值}
设 $f$ 为周期 $2$ 的函数，在 $[-1,1)$ 上为
\[f(x)=\begin{cases}1-3x,&-1\le x<0,\\x^2+1,&0\le x<1.\end{cases}\]
记其傅里叶级数的和函数为 $S$。不计算全部傅里叶系数，求 $S(-7),S(0),S(\frac12)$，并写出所有满足 $S(x)\ne f(x)$ 的点。说明所用收敛定理的条件。
\end{exercise}
\begin{exercise}{785}{余弦延拓与跳跃}
设 $f(x)=2x$（$0\le x\le\frac12$），$f(x)=x^2+1$（$\frac12<x\le1$），并令
\[a_n=2\int_0^1f(x)\cos(n\pi x)\,\dd x\quad(n=0,1,\ldots),\qquad
S(x)=\frac{a_0}2+\sum_{n=1}^{\infty}a_n\cos(n\pi x).\]
说明该余弦级数对应的偶延拓与周期，求 $S(-\frac94),S(\frac12),S(1)$，并求常数项 $a_0/2$。
\end{exercise}
\begin{exercise}{786}{正弦延拓的边界值}
设 $f(x)=2x^2$（$0\le x<1$），$b_n=2\int_0^1f(x)\sin(n\pi x)\,\dd x$，$S(x)=\sum_{n=1}^{\infty}b_n\sin(n\pi x)$。说明相应奇延拓的周期与跳跃点，求 $S(-\frac12)+S(9)+S(10)$；再求 $b_n$ 并用它验证所有整数点处的和。
\end{exercise}
\begin{exercise}{787}{傅里叶恒等式的取点求和}
已知在 $0\le x\le2$ 上成立
\[(x-1)^2=\frac13+\frac4{\pi^2}\sum_{n=1}^{\infty}\frac{\cos(n\pi x)}{n^2}.\]
选择适当的 $x$ 求 $\sum n^{-2}$ 与 $\sum(-1)^{n-1}n^{-2}$，再求 $\sum_{n=1}^{\infty}(2n-1)^{-2}$。每次代入均须写明所用的点及级数收敛的理由。
\end{exercise}
\begin{exercise}{788}{余弦展开与周期端点}
将 $f(x)=x^2-x$（$0\le x\le1$）展开成周期为 $2$ 的余弦级数，写出全部傅里叶系数及在端点处的和；利用 $x=\frac12$ 处的展开求 $\sum_{m=1}^{\infty}\frac{(-1)^m}{m^2}$。
\end{exercise}
\begin{exercise}{789}{单点赋值与傅里叶和}
设 $f$ 为周期 $2$ 的函数，在 $(-1,0)$ 上等于 $-1$，在 $(0,1)$ 上等于 $1$，并令 $f(0)=c,f(1)=f(-1)=d$。证明傅里叶系数与 $c,d$ 无关，求傅里叶级数在全部整数点处的和，并确定使该和等于函数值的 $c,d$。
\end{exercise}
\begin{exercise}{790}{余弦展开与参数级数恒等式}
将 $f_a(x)=x^2+ax$（$0\le x\le1$）作偶延拓后再作周期 $2$ 延拓，求其余弦级数。确定使所有奇数阶余弦系数消失的 $a$，再利用 $x=0,1/2$ 的取值求出相应的两个数项级数恒等式，明确端点收敛依据。
\end{exercise}
\subsubsection{综合提高题}
\begin{exercise}{791}{通项趋零与复合极限}
已知 $\sum_{n=1}^{\infty}(u_n-\frac{3n}{n+1})$ 收敛。求 $u_n$ 的极限及 $\lim_{n\to\infty}(u_n^2-2u_n+3)$；判断能否据此断言 $\sum(u_n-3)$ 收敛，并证明你的结论。
\end{exercise}
\begin{exercise}{792}{错位级数与首项修正}
已知 $\sum_{k=1}^{n}u_k=\frac{n}{2n+1}$。求 $u_n$ 及 $\sum_{n=1}^{\infty}(u_n+u_{n+1}-u_{n+2})$ 的和；进一步求其前 $N$ 项和与总和之差，说明错位相加时哪些首项需要单独保留。
\end{exercise}
\begin{exercise}{793}{三重分母裂项与精确余项}
按级数收敛定义求 $\sum_{n=1}^{\infty}\frac1{n(n+1)(n+2)}$ 的和，写出前 $N$ 项和及精确余项；证明余项是 $N^{-2}$ 阶，并求 $N^2R_N$ 的极限。
\end{exercise}
\begin{exercise}{794}{幂和恒等式与级数求和}
先用数学归纳法证明 $1^3+2^3+\cdots+n^3=(n(n+1)/2)^2$，再求 $\sum_{n=1}^{\infty}\frac1{\sqrt{1^3+2^3+\cdots+n^3}}$ 的和；写出用前 $N$ 项近似时的精确误差。
\end{exercise}
\begin{exercise}{795}{通项收敛与部分和有界的区别}
设 $u_n>0$，$S_n=\sum_{k=1}^{n}u_k$。证明 $\{S_n\}$ 有界是 $\{u_n\}$ 收敛的充分条件，说明该条件是否必要并给出反例；再证明 $\{S_n\}$ 有界与正项级数 $\sum u_n$ 收敛等价，区分两种收敛对象。
\end{exercise}
\begin{exercise}{796}{单调部分和的倒数比较}
设 $u_n>0$，$S_n=\sum_{k=1}^{n}u_k\to\infty$。证明对 $n\ge2$ 有 $\frac{u_n}{S_n^2}\le\frac1{S_{n-1}}-\frac1{S_n}$，据此证明 $\sum\frac{u_n}{S_n^2}$ 收敛并给出关于 $u_1$ 的上界；讨论证明是否实际需要原级数发散。
\end{exercise}
\begin{exercise}{797}{抵消精度与参数辨析}
讨论 $\sum_{n=1}^{\infty}(\frac1{\sqrt n+1}-\frac1{\sqrt n}+\frac cn)$ 随 $c\in\mathbb R$ 的敛散性，给出 $c=1$ 时的收敛类型；在收敛情形下将通项化为单个正分式并求其与 $n^{-3/2}$ 的比值极限。
\end{exercise}
\begin{exercise}{798}{积分定义通项与尾和控制}
设 $u_n=\int_n^{n+1}\frac{e^{-x}}{1+\sin^2n}\,\dd x$。证明 $\sum u_n$ 收敛，给出级数和的两个明确常数界及前 $N$ 项余项的指数上界；说明分母中的振荡为何不改变收敛性。
\end{exercise}
\begin{exercise}{799}{一阶与二阶无穷小的级数差别}
分别判定 $\sum(e^{1/n}-1)$、$\sum\ln(1+1/n)$、$\sum\sin(1/n)$、$\sum(1-\cos(1/n))$ 的敛散性，写出对应等价无穷小；再求使 $\sum(1-\cos(1/n))^p$ 收敛的全部 $p>0$。
\end{exercise}
\begin{exercise}{800}{极限指数与幂次阈值}
先判定 $\sum_{n=1}^{\infty}\frac{(n+2)^n}{n^{n+2}}$，再对 $c>-1,p\in\mathbb R$ 讨论 $\sum_{n=1}^{\infty}\frac{(n+c)^n}{n^{n+p}}$ 的敛散性；写出与 $p$ 级数的比较极限，单独解释临界值。
\end{exercise}
\begin{exercise}{801}{交错分母产生的非交错余项}
判定 $\sum_{n=1}^{\infty}\frac{(-1)^{n+1}}{\sqrt n+(-1)^{n+1}}$ 的敛散性。对 $n\ge2$ 有理化分母，将通项分为一个交错部分与一个同号部分，分别判定；若发散，说明部分和趋于何处，不能仅凭符号交替作答。
\end{exercise}
\begin{exercise}{802}{参数分母与隐藏调和项}
设 $k>1$。把 $\frac{(-1)^n}{\sqrt{kn}-(-1)^n}$ 有理化，判定相应级数的敛散性并说明部分和的极限行为；解释参数 $k$ 是否改变结论，以及 $k>1$ 如何保证全部通项有定义。
\end{exercise}
\begin{exercise}{803}{整数相位与小量振荡}
设 $k>0$，判定 $\sum_{n=1}^{\infty}\sin(\pi\sqrt{n^2+k})$ 的收敛类型。令 $\delta_n=\pi(\sqrt{n^2+k}-n)$，证明 $\delta_n$ 单调趋于零并求 $n\delta_n$ 的极限，再判断 $k=0$ 时的情形。
\end{exercise}
\begin{exercise}{804}{奇偶子级数的联立关系}
设 $u_n>0$，$\sum u_n$ 发散而 $\sum(-1)^{n-1}u_n$ 收敛。证明奇数项子级数与偶数项子级数都发散，判定 $\sum(u_{2n-1}+u_{2n})$ 与 $\sum(u_{2n-1}-u_{2n})$ 的敛散性；用各自前 $N$ 项和说明所有结论。
\end{exercise}
\begin{exercise}{805}{分式变量代换与收敛域}
求 $\sum_{n=1}^{\infty}\frac1{2n+1}(\frac{1-x}{1+x})^n$ 的实收敛域，明确排除无定义点并检验全部边界。对任意 $0<a<b<\infty$，给出 $[a,b]$ 上的几何尾和界，说明其和函数在 $(0,\infty)$ 上连续。
\end{exercise}
\begin{exercise}{806}{收敛域的平移反推}
已知 $\sum_{n=1}^{\infty}\frac{(-1)^{n-1}}n(x-a)^n$ 对每个 $x>0$ 都发散，而在 $x=0$ 收敛。求 $a$，写出完整收敛域及和函数，并说明 $x=0$ 处是绝对收敛还是条件收敛。
\end{exercise}
\begin{exercise}{807}{单调系数与端点信息}
设实数列 $a_n$ 单调递减趋于零，且 $S_n=\sum_{k=1}^{n}a_k$ 无界。求 $\sum_{n=1}^{\infty}a_n(x-2)^n$ 的收敛半径与收敛域；证明各 $a_n$ 非负，分别解释内点、左端点和右端点的收敛类型。
\end{exercise}
\begin{exercise}{808}{二次分母拆分求和}
求 $S(x)=\sum_{n=2}^{\infty}\frac{x^n}{n^2-1}$ 的完整收敛域及和函数，单独写明 $x=0$ 的连续取值；求 $S(1/2)$，并由端点级数或极限求 $S(1)$。
\end{exercise}
\begin{exercise}{809}{阶乘加权和与可去点}
求 $S(x)=\sum_{n=1}^{\infty}\frac n{(n+1)!}x^n$ 的收敛域及和函数，明确 $x=0$ 的取值；求 $S(1)$，再由展开系数求 $S'(0)$，核验所得分式表达式在原点的可去性。
\end{exercise}
\begin{exercise}{810}{微分方程的多项式解与求和}
对每个正整数 $n$，求满足 $xy_n'-2ny_n=0$ 且 $y_n(1)=\frac1{n(2n-1)}$ 的多项式 $y_n(x)$，并证明在多项式类中的唯一性。求 $\sum_{n=1}^{\infty}(-1)^{n-1}y_n(x)$ 的收敛域及和函数，分别处理 $x=0$ 和两个端点。
\end{exercise}
\begin{exercise}{811}{反常面积与对数求和}
设 $D_n$ 为 $x\ge e^3,0\le y\le\frac1{x(\ln x)^{n+1}}$ 的无界区域，$n\ge1$。求其面积 $A_n$，证明 $\sum(-1)^{n-1}A_n$ 收敛并求和；再将 $e^3$ 改为 $e^b$（$b>0$），讨论所得面积级数随 $b$ 的敛散性。
\end{exercise}
\begin{exercise}{812}{两个对数的平移展开}
利用 $\ln(1+t)$ 的展开将 $\ln\frac{x}{1+x}$ 展开成 $(x-1)$ 的幂级数，写出通项系数与收敛域；单独判断两个端点，并说明它们是否属于原函数的定义域。
\end{exercise}
\begin{exercise}{813}{反正切恒等式的分支与端点}
设 $f(x)=x\arctan\frac{x+2}{x-2}$，其中 $\arctan$ 取值于 $(-\pi/2,\pi/2)$。先在 $-2<x<2$ 推导适用的反正切恒等式，再展开成 $x$ 的幂级数；分别求级数收敛域及展开等于原函数的范围，特别解释 $x=2$ 处级数收敛与原函数无定义的关系。
\end{exercise}
\begin{exercise}{814}{乘积展开的有限阶系数}
设 $f(x)=xe^{-x^2}\sin2x$。不反复求导，求 $f^{(4)}(0)$ 与 $f^{(6)}(0)$；说明提取这两阶系数各需将两个因子展开到什么阶，并证明全部奇数阶原点导数为零。
\end{exercise}
\begin{exercise}{815}{柯西乘积与一般导数}
设 $f(x)=(\sum_{n=1}^{\infty}x^n)^2$，$|x|<1$。通过柯西乘积求 $x^n$ 的系数及 $f^{(n)}(0)$（$n\ge2$），写出 $n=0,1$ 时的导数值；再用有理函数形式验证，并说明级数乘法合法的依据。
\end{exercise}
\begin{exercise}{816}{方波的跳跃平均与系数}
设 $f(x)=\operatorname{sgn}(\cos x)$，并规定 $\operatorname{sgn}(0)=0$。求其傅里叶级数和函数 $S$，计算 $S(-\pi/2)+S(3\pi)$；利用偶对称性求全部傅里叶系数，说明在跳跃点的取值依据。
\end{exercise}
\begin{exercise}{817}{余弦系数的奇偶分离}
将 $f(x)=1-x$（$0\le x\le1$）作偶延拓并以 $2$ 为周期延拓，记余弦系数为 $a_n$。求全部 $a_n$，求 $\sum_{n=1}^{\infty}a_{2n}$ 与 $\sum_{n=1}^{\infty}a_{2n-1}$；再用 $x=0,1$ 处的傅里叶和核验这两个数项和。
\end{exercise}
\begin{exercise}{818}{分段周期函数与系数总和}
设 $f$ 为周期 $2$ 的函数，在 $-1<x\le0$ 上 $f(x)=x^2$，在 $0<x\le1$ 上 $f(x)=x-1$。令 $a_n=\int_{-1}^{1}f(x)\cos(n\pi x)\,\dd x$（$n\ge0$）。求傅里叶级数的和函数在 $[-1,1]$ 上的完整表达式，并求 $\sum_{n=0}^{\infty}a_n$ 与 $\sum_{n=0}^{\infty}(-1)^na_n$；说明含 $a_0$ 的数项和与傅里叶常数项 $a_0/2$ 之间的差别。
\end{exercise}
